\documentclass[10pt,a4paper]{amsart}
\usepackage[margin=19mm]{geometry}
\usepackage{amsmath,amssymb,mathtools}
\usepackage[scr=rsfs,cal=euler]{mathalfa}
\usepackage{xfrac}
\usepackage[unicode,hidelinks,bookmarks=false]{hyperref}
\newcommand{\ie}{i.e.\ }
\newcommand{\cf}{cf.\ }
\newcommand{\resp}{respectively\ }
\newcommand{\Subsection}{\subsection}
\providecommand{\nobreakdash}{\nobreak\hbox{-}}
\setlength{\emergencystretch}{2em}
\multlinegap0pt
\usepackage{amssymb}
\usepackage{dsfont}
\usepackage{tikz}
\newtheorem{theorem}{Theorem}
\title{Foundations and Classification of Invariant Subalgebras of Grassmann Algebra}
\author{Mithat Konuralp Demir}
\date{Source: arXiv:2603.09766v1 (CC BY 4.0)}
\begin{document}
\maketitle

\noindent\emph{Source and licence.} Foundations and Classification of Invariant Subalgebras of Grassmann Algebra. Version arXiv:2603.09766v1 (CC BY 4.0), distributed by arXiv under the Creative Commons Attribution 4.0 International licence, http://creativecommons.org/licenses/by/4.0/, as recorded in the arXiv licence metadata for this exact version and shown on the abstract page https://arxiv.org/abs/2603.09766v1. Full licence notice: \texttt{licenses/arxiv-2603.09766-CC-BY-4.0.txt}.\par\smallskip

\noindent\textit{Editorial scope.} Counted unit: the article's theorem theorem (source0.tex lines 841--844). The article's complete original proof of it is delivered with it (source0.tex lines 845--879, one \texttt{proof} environment of 35 lines). No mathematics is written, repaired or re-derived: the statement and the proof are the article's own lines, in the article's own notation, and no retained line is substituted or re-flowed.  No further source-local context is needed: every cross-reference of every retained line resolves inside the two delivered spans.  Nothing was omitted from any retained span, so the package carries no omission record.  No retained line contains a citation, so the wrapper carries no bibliography and no bibliography entry.  No retained line contains a figure environment, an image-inclusion command or drawing code, so no image is delivered and the package declares no asset dependency.  Editorial formatting only, in addition: an AMS \texttt{amsart} preamble that restates the article's own declarations and macros used by the retained lines, namely the environments \texttt{theorem} on separate counters, together with the standard packages those lines need.  The retained environments print in this document as Theorem 1 in order of appearance, whereas the article prints the same items with the numbering of its own full document.  The complete native source of the article as distributed in the arXiv e-print for this exact version is bundled unchanged as \texttt{sources/196084/source0.tex}.
\begin{theorem}
    Leibniz formula and Cofactor expansion are equivalent:
    $$det(A) = \sum_{\sigma \in S_n}(sgn \,\sigma)\,\prod_{i = 0}^na_{i, \sigma(j)} = S_i = \sum_{j = 1}^n a_{ij}\,C{ij}$$ 
\end{theorem}

\begin{proof}
    We prove the equality of formulas by induction on $n$. Also, we consider, without loss of generalization, $1^{st}$ row cofactor expansion $S_1$.
    
    Base case: $n = 1$
    We consider $A = [a_{11}]$
    Sign of the identity permutation is $+1$ and cofactor in this case is defined to be $1$. Thus base case holds.

    Inductive hypothesis:
    Assume for all $(n-1)\times(n-1)$ matrices, the Leibniz formula equals cofactor expansion, i.e. the determinant.

    Let $\sigma \in S_n$ and $\bar{\sigma}$ be the restriction of $\sigma$ to $\{2, ..., n\}$. $\sigma(1)$ cannot be in the image of $\bar{\sigma}$ and since there is $n-1$ elements in the image and we have $n$ possible slots, by pigeonhole principle we have one option, which we denote as $\sigma(1) := j_0$ 
    Then for the particular extension $\sigma$, we have:
    $$
    \prod_{i = 1}^na_{i\sigma(i)} = a_{1j_0}\prod_{i = 2}^n a_{i\bar{\sigma}(i)}
    $$

    Moving $j_0$ to the first position, such that $\sigma(1) = j_0$ contributes $(-1)^{1+j_0}$ to the sign relative to $\bar{\sigma}$, the permutation of the remaining indices. So we relate signs of $\sigma$ and $\bar{\sigma}$ as follows:
    $$(sgn\;\sigma) = (-1)^{1+j_0}(sgn\;\bar{\sigma})$$ 

    By inductive hypothesis, the determinant of the minor $M_{1j_0}$, which is a $(n - 1) \times (n - 1)$ matrix is:
    $$det(M_{1j_0}) = \sum_{\bar{\sigma} \in S_{n-1}}(sgn\;\bar{\sigma})\prod_{i = 2}^n a_{i \bar{\sigma}(i)}$$

    For each $\sigma$ with $\sigma(1) = j_0$:
    $$
    (sgn\;\sigma)\prod_{i=1}^n a_{i\sigma(i)} = (-1)^{1+j_0}(sgn\;\bar{\sigma}) a_{1j_0} \prod_{i = 2}^n a_{i \bar{\sigma}(i)} = a_{1j_0} (-1)^{1+j_0}det(M_{1j_0})
    $$

    Now we partition $S_n$ according to $\sigma(1) = j$ for $j \in \{1, ..., n\}$. Then summing over all $\sigma$:

    \begin{align*}
        det(A) &= \sum_{\sigma \in S_n}(sgn\;\sigma)\prod_{i = 1}^n a_{i\sigma(i)}\\
        &= \sum_{j = 1}^n a_{1j}(-1)^{1+j} \sum_{\bar{\sigma} \in S_{n-1}} \prod_{i = 2}^n a_{i\bar{\sigma}(i)}\\
        &= \sum_{j = 1}^n a_{1j}(-1)^{1+j}\,det(A_{1j}) = \sum_{j = 1}^n a_{1j}\,C_{1j} 
    \end{align*}
\end{proof}

\end{document}
